\section{Related Work}
\label{sec:related_work}

\para{Recursive self-training and model collapse.}
The central observation from the model collapse literature is that recursive training on generated data can erase distributional support before broad task quality visibly fails. Shumailov et al.\ show that repeated training on generated data drives tail loss and distribution drift across multiple modalities \cite{shumailov2023curse}. Alemohammad et al.\ describe related self-consuming dynamics and emphasize the role of fresh real data in preventing autophagy \cite{alemohammad2023self}. For language specifically, Briesch et al.\ find that diversity deteriorates long before correctness collapses in a self-consuming loop \cite{briesch2023large}. We build directly on this diversity-first framing, but we ask a different question: whether collapse behavior changes when the recursive loop is treated as a population rather than a single lineage.

\para{Accumulation and correction as anti-collapse mechanisms.}
Several recent papers argue that collapse is not inevitable once the data regime changes. Gillman et al.\ show that corrective operators can stabilize recursive generative training \cite{gillman2024selfcorrecting}. Gerstgrasser et al.\ give a direct counterexample to inevitability by showing that accumulation of real and synthetic data can remain stable when replacement fails \cite{gerstgrasser2024model}. These papers motivate one of our core design choices. We treat replacement and accumulation as first-order experimental factors, because any claim about a minimum viable population is confounded if human refresh is ignored. Our results support this emphasis on data governance, although in our bounded setting accumulation helps only modestly on average.

\para{Population diversity and operational individuality.}
The minimum viable population idea comes from biology and ecology, but the closest operational precedent for our setting comes from multi-agent learning. McKee et al.\ show that population and environment diversity can affect held-out generalization, and they measure behavioral diversity rather than assuming it from nominal counts \cite{mckee2021quantifying}. That distinction is crucial for LLM ecosystems. Agent-society and role-playing work such as CAMEL, Generative Agents, and RoleLLM provides mechanisms for constructing persistent or role-conditioned differences, but these papers do not study recursive collapse itself \cite{li2023camel,park2023generative,wang2023rolellm}. We borrow their intuition that roles and memory can define agents behaviorally, then test whether that distinctness is enough to protect a recursive ecosystem.

\para{Positioning our work.}
Our paper sits between recursive collapse studies and multi-agent diversity studies. Unlike prior single-lineage collapse work, we vary population size and prompt-level distinctness explicitly. Unlike agent-society work, we focus on recursive inheritance stability rather than task coordination. The result is a population-level proxy experiment that offers two concrete outputs: an operational answer to what counts as an individual, and direct evidence that prompt-level distinctness alone does not create a simple minimum viable population threshold.
